\section{Introduction}
\label{sec:introduction}

Modern scientific simulations and instruments generate massive volumes of multidimensional floating-point data. The resulting files can be large individually, and complete datasets may comprise multiple physical fields and time steps. Managing these growing collections calls for effective data reduction that preserves the fidelity required by scientific users~\cite{cappello2019use,di2025survey,cappello2025support}.

Error-bounded lossy compression provides such data reduction by substantially reducing data volume while guaranteeing that reconstruction errors satisfy a user-specified bound. Representative compressors include SZ3~\cite{zhao2021optimizing,liang2022sz3}, QoZ~\cite{liu2022dynamic,liu2024high}, ZFP~\cite{lindstrom2014fixed}, and SPERR~\cite{li2023lossy}. Among them, prediction-based compressors such as SZ3 and QoZ support simple predictors that estimate data values through weighted combinations of neighboring reconstructed values. These compressors combine high compression ratios with high throughput on many scientific datasets.

Prediction-based compression benefits from configurations adapted to the data, a principle also explored in data-adaptive storage encoding for database systems~\cite{kuschewski2023btrblocks,afroozeh2023alp}. State-of-the-art compressors, including SZ3 and QoZ, tune prediction stencils (the relative locations of reference points), interpolation coefficients, and prediction directions for different datasets and error bounds. This tuning helps them identify well-performing configurations within their search spaces and improve rate--distortion performance, achieving smaller compressed sizes at comparable reconstruction quality.

Nevertheless, room for improvement remains. A larger search space offers more opportunities to find better configurations, but realizing these gains requires accurate candidate evaluation. Both generally increase tuning cost, motivating a more efficient selection mechanism.

We address this challenge with ASZ, a prediction-based error-bounded lossy compressor. ASZ combines a broader choice of predictor and encoding configurations with accurate, low-cost tuning. Our main contributions are fourfold:

\begin{itemize}
\item We introduce tunable-coefficient interpolation predictors to expand the coefficient search space and reassess which predictors remain competitive.
\item We combine an accurate rate--distortion proxy with reusable coefficient-independent statistics to make coefficient selection more efficient.
\item We identify quantization-bin serialization order as an independent tuning dimension that can improve lossless compression without altering reconstruction.
\item We develop reverse-greedy search to select layer-wise prediction directions and quantization-bin serialization orders efficiently. The selected configurations are stored with the selected predictor in an offline plan that can be reused across related files.
\end{itemize}

We evaluate ASZ on 18 three-dimensional scientific fields drawn from simulations and experimental instruments. At equal PSNR, ASZ achieves $2.51\times$ and $2.03\times$ the compression ratios of SZ3 and QoZ, respectively, on average. Its end-to-end runtime, including preprocessing, compression, and decompression, remains comparable to QoZ. Experiments on multi-file datasets also show that reusing configurations across related files preserves the compression advantage and further amortizes offline tuning cost.

The remainder of this paper is organized as follows. Section~\ref{sec:related-work} reviews related work. Section~\ref{sec:background} formulates our compression and tuning objectives, then introduces the shared SZ compression pipeline and reviews interpolation and tuning in SZ3 and QoZ. Section~\ref{sec:configuration-space} examines how ASZ expands and prunes the configuration space and develops accurate, efficient coefficient evaluation. Section~\ref{sec:asz} presents the joint search and the complete ASZ compressor. Section~\ref{sec:evaluation} evaluates rate--distortion performance, runtime, robustness, and configuration reuse across files. The remaining sections discuss limitations and conclude the paper.
